\section{AMC 8 2000}
        \begin{enumerate}
        	\item (\textit{AMC 8 2000}) Aunt Anna is $42$ years old. Caitlin is $5$ years younger than Brianna, and Brianna is half as old as Aunt Anna. How old is Caitlin?


 $\mathrm{(A)}\ 15\qquad\mathrm{(B)}\ 16\qquad\mathrm{(C)}\ 17\qquad\mathrm{(D)}\ 21\qquad\mathrm{(E)}\ 37$


 
 \textbf{Answer: }B



\textbf{Solution}

Brianna is $\frac{1}{2} \cdot 42 = 21$ years old, so Caitlin is $21-5=\boxed{\textbf{(B)}\ 16}$ years old.

 


	\item (\textit{AMC 8 2000}) Which of these numbers is less than its reciprocal?


 $\text{(A)}\ -2 \qquad \text{(B)}\ -1 \qquad \text{(C)}\ 0 \qquad \text{(D)}\ 1 \qquad \text{(E)}\ 2$


 
 \textbf{Answer: }A



\textbf{Solution 1}

The reciprocal of $-2$ is $-\frac{1}{2}$.  Since $-2 < -\frac{1}{2}$, $\boxed{\textbf{(A)}\ {-2}}$ is correct.

 The reciprocal of $-1$ is $-1$.  Since $-1 = -1$, (B) is incorrect.

 The reciprocal of $0$ is undefined, so (C) is incorrect.

 The reciprocal of $1$ is $1$.  Since $1 = 1$, (D) is incorrect.

 The reciprocal of $2$ is $\frac{1}{2}$.  Since $2 > \frac{1}{2}$, (E) is incorrect.

 



\textbf{Solution 2}

Let $x$ be a real number that is less than its reciprocal, that is, $x < \frac{1}{x}$.

 Since $\frac{1}{x}$ must be defined, $x \neq 0$.

 If $x$ is positive, then $x^2 < 1$, so $0 < x < 1$.

 If $x$ is negative, then $x^2 > 1$, so $x < -1$.

 Therefore, $x < \frac{1}{x}$ if and only if either $x < -1$ or $0 < x < 1$.  Of the answer choices, only $\boxed {\textbf{(A)}\ {-2}}$ is a possible value of $x$.

 



\textbf{Solution 3}

Starting again with $x < \frac{1}{x}$, we avoid multiplication by $x$.   Instead, move everything to the left, and find a common denominator:

 $x < \frac{1}{x}$

 $x - \frac{1}{x} < 0$

 $\frac{x^2 - 1}{x} < 0$

 $\frac{(x+1)(x-1)}{x} < 0$

 Divide this expression at $x=-1$, $x=0$, and $x=1$, as those are the three points where the expression on the left will "change sign".

 If $x<-1$, all three of those terms will be negative, and the inequality is true.  Therefore, $(-\infty, -1)$ is part of our solution set.

 If $-1 < x < 0$, the $(x+1)$ term will become positive, but the other two terms remain negative.  Thus, there are no solutions in this region.

 If $0 < x < 1$, then both $(x+1)$ and $x$ are positive, while $(x-1)$ remains negative.  Thus, the entire region $(0, 1)$ is part of the solution set.

 If $1 < x$, then all three terms are positive, and there are no solutions.

 At all three "boundary points", the function is either $0$ or undefined.  Therefore, the entire solution set is $(-\infty, -1) \cup (0, 1)$, and the only option in that region is $x=-2$, leading to $\boxed{A}$.

 


	\item (\textit{AMC 8 2000}) How many whole numbers lie in the interval between $\frac{5}{3}$ and $2\pi?$


 $\text{(A)}\ 2 \qquad \text{(B)}\ 3 \qquad \text{(C)}\ 4 \qquad \text{(D)}\ 5 \qquad \text{(E)}\ \text{infinitely many}$


 
 \textbf{Answer: }D



\textbf{Solution}

Note that $1 < \frac{5}{3} < 2$ while $6 < 2\pi < 7$.  Therefore, the interval from $\frac{5}{3}$ to $2\pi$ contains $\boxed{\textbf{(D)}\ 5}$ whole numbers: $2$, $3$, $4$, $5$, and $6$.

 


	\item (\textit{AMC 8 2000}) In $1960$ only $5\%$ of the working adults in Carlin City worked at home. By $1970$ the "at-home" work force increased to $8\%$. In $1980$ there were approximately $15\%$ working at home, and in $1990$ there were $30\%$. The graph that best illustrates this is


 
                    \begin{figure}[H]
                        \centering
                        \includegraphics[width=0.5\textwidth]{Problem/AMC_8_AMC_2000_4_Problem_Figure_1.png}
                    \end{figure}
                    


 
 \textbf{Answer: }E



\textbf{Solution 1}

The data are $1960 (5\%)$, $1970 (8\%)$, $1980 (15\%)$, and $1990 (30\%)$. Only one of these graphs has the answer and that is choice $\boxed{E}$.

 



\textbf{Solution 2}

We can look at the graphs and note that the only one that has all the correct points is $\boxed{E}$

 


	\item (\textit{AMC 8 2000}) Each principal of Lincoln High School serves exactly one $3$-year term. What is the maximum number of principals this school could have during an $8$-year period?


 $\text{(A)}\ 2 \qquad \text{(B)}\ 3 \qquad \text{(C)}\ 4 \qquad \text{(D)}\ 5 \qquad \text{(E)}\ 8$


 
 \textbf{Answer: }C



\textbf{Solution 1}

Suppose the current principal's term will end in $1$ year.  Of the next $8$ years, the current principal will serve the first year; the next principal will serve the second, third, and fourth years; a third principal will serve the fifth, sixth, and seventh years; and a fourth principal will serve the eighth year.  Therefore, $4$ different principals will serve in the next $8$ years.

 Consider a period that includes at least $5$ different principals.  During this period, all principals except the first and last must serve complete terms. So at least $3$ principals serve complete terms, and thus this period is at least $9$ years long.  Therefore, an $8$-year period may not have $5$ or more different principals.

 Thus there can be at most $\boxed{\textbf{(C)}\ 4}$ principals in an $8$-year period.

 



\textbf{Solution 2}

In an $8$-year period, at most $\left\lfloor{\frac{8}{3}}\right\rfloor = 2$ principals could serve complete terms.  Furthermore, at most $2$ principals could serve incomplete terms during this period: those who are in office at the start and end of this period.  Therefore, at most $2 + 2 = \boxed{\textbf{(C)}\ 4}$ principals could be active during this period (which would happen if the acting principal at the start of the period has $1$ year left of their term).

 


	\item (\textit{AMC 8 2000}) Figure $ABCD$ is a square. Inside this square three smaller squares are drawn with the side lengths as labeled. The area of the shaded $L$-shaped region is


 
                    \begin{figure}[H]
                        \centering
                        \includegraphics[width=0.5\textwidth]{Problem/AMC_8_AMC_2000_6_Problem_Figure_1.png}
                    \end{figure}
                    


 $\text{(A)}\ 7 \qquad \text{(B)}\ 10 \qquad \text{(C)}\ 12.5 \qquad \text{(D)}\ 14 \qquad \text{(E)}\ 15$


 
 \textbf{Answer: }A



\textbf{Solution 1}

The side of the large square is $1 + 3 + 1 = 5$, so the area of the large square is $5^2 = 25$.

 The area of the middle square is $3^2$, and the sum of the areas of the two smaller squares is $2 * 1^2 = 2$.

 Thus, the big square minus the three smaller squares is $25 - 9 - 2 = 14$.  This is the area of the two congruent L-shaped regions.

 So the area of one L-shaped region is $\frac{14}{2} = 7$, and the answer is $\boxed{A}$

 



\textbf{Solution 2}

The shaded area can be divided into two regions:  one rectangle that is 1 by 3, and one rectangle that is 4 by 1.  (Or the reverse, depending on which rectangle the 1 by 1 square is "joined" to.)  Either way, the total area of these two regions is $3 + 4 = 7$, and the answer is $\boxed{A}$.

 



\textbf{Solution 3}

Chop the entire 5 by 5 region into $25$ squares like a piece of graph paper.  When you draw all the lines, you can count that only $7$ of the small 1 by 1 squares will be shaded, giving $\boxed{A}$ as the answer.

 



\textbf{Solution 4}

In the bottom left corner of the 5 by 5 square there is a 4 by 4 square which has an area of $4\cdot4=16$. In the top right of that 4 by 4 square is a 3 by 3 square with an area of $3\cdot3=9$. When we remove the 3 by 3 square from the 4 by 4 square we get the L-shaped figure so our answer is $16-9=\boxed{\text{(A)}\ 7}$

 


	\item (\textit{AMC 8 2000}) What is the minimum possible product of three different numbers of the set $\{-8,-6,-4,0,3,5,7\}$?


 $\text{(A)}\ -336 \qquad \text{(B)}\ -280 \qquad \text{(C)}\ -210 \qquad \text{(D)}\ -192 \qquad \text{(E)}\ 0$


 
 \textbf{Answer: }B



\textbf{Solution}

If the product of $3$ selected numbers is negative, either all $3$ numbers are negative or $1$ number is negative and $2$ numbers are positive.

 The minimum product of $3$ negative numbers from this set is $(-8)(-6)(-4) = -192$, while the minimum product of $1$ negative number and $2$ positive numbers from this set is $(-8)(5)(7) = -280$.

 Therefore, the minimum product of $3$ numbers from this set is $\boxed{\textbf{(B)}\ {-280}}$.

 


	\item (\textit{AMC 8 2000}) Three dice with faces numbered 1 through 6 are stacked as shown. Seven of the eighteen faces are visible, leaving eleven faces hidden (back, bottom, between). The total number of dots NOT visible in this view is


 
                    \begin{figure}[H]
                        \centering
                        \includegraphics[width=0.5\textwidth]{Problem/AMC_8_AMC_2000_8_Problem_Figure_1.png}
                    \end{figure}
                    


 $\text{(A)}\ 21 \qquad \text{(B)}\ 22 \qquad \text{(C)}\ 31 \qquad \text{(D)}\ 41 \qquad \text{(E)}\ 53$


 
 \textbf{Answer: }D



\textbf{Solution}

The numbers on one die total $1+2+3+4+5+6 = 21$, so the numberson the three dice total $63$. Numbers $1, 1, 2, 3, 4, 5, 6$ are visible, and these total $22$.This leaves $63 - 22 = \boxed{\text{(D) 41}}$ not seen.

 


	\item (\textit{AMC 8 2000}) Three-digit powers of $2$ and $5$ are used in this "cross-number" puzzle. What is the only possible digit for the outlined square?

 \[\begin{array}{lcl} \textbf{ACROSS} & & \textbf{DOWN} \\ \textbf{2}.~ 2^m & & \textbf{1}.~ 5^n \end{array}\]


 
                    \begin{figure}[H]
                        \centering
                        \includegraphics[width=0.5\textwidth]{Problem/AMC_8_AMC_2000_9_Problem_Figure_1.png}
                    \end{figure}
                    


 $\text{(A)}\ 0 \qquad \text{(B)}\ 2 \qquad \text{(C)}\ 4 \qquad \text{(D)}\ 6 \qquad \text{(E)}\ 8$


 
 \textbf{Answer: }D



\textbf{Solution}

The $3$-digit powers of $5$ are $125$ and $625$, so space $2$ is filled with a $2$.The only $3$-digit power of $2$ beginning with $2$ is $256$, so the outlined block is filled witha $\boxed{\text{(D) 6}}$.

 


	\item (\textit{AMC 8 2000}) Ara and Shea were once the same height. Since then Shea has grown 20\% while Ara has grown half as many inches as Shea. Shea is now 60 inches tall. How tall, in inches, is Ara now?


 $\text{(A)}\ 48 \qquad \text{(B)}\ 51 \qquad \text{(C)}\ 52 \qquad \text{(D)}\ 54 \qquad \text{(E)}\ 55$


 
 \textbf{Answer: }E



\textbf{Solution}

Shea has grown $20\%$, if x was her original height, then $1.2x = 60$, so she was originally $\frac{60}{1.2}=50$ inches tall which is a $60 - 50 = 10$ inch increase.  Ara also started off at $50$ inches.  Since Ara grew half as much as Shea, Ara grew $\frac{10}{2} = 5$ inches.  Therefore, Ara is now $50+5=55$ inches tall which is choice $\boxed{E}.$

 


	\item (\textit{AMC 8 2000}) The number $64$ has the property that it is divisible by its units digit. How many whole numbers between 10 and 50 have this property?


 $\text{(A)}\ 15 \qquad \text{(B)}\ 16 \qquad \text{(C)}\ 17 \qquad \text{(D)}\ 18 \qquad \text{(E)}\ 20$


 
 \textbf{Answer: }C



\textbf{Solution}

Casework by the units digit $u$ will help organize the answer.

 $u=0$ gives no solutions, since no real numbers are divisible by $0$

 $u=1$ has $4$ solutions, since all numbers are divisible by $1$.

 $u=2$ has $4$ solutions, since every number ending in $2$ is even (ie divisible by $2$).

 $u=3$ has $1$ solution: $33$.  $\pm 10$ or $\pm 20$ will retain the units digit, but will stop the number from being divisible by $3$.  $\pm 30$ is the smallest multiple of $10$ that will keep the number divisible by $3$, but those numbers are $3$ and $63$, which are out of the range of the problem.

 $u=4$ has $2$ solutions:  $24$ and $44$.  Adding or subtracting $10$ will kill divisibility by $4$, since $10$ is not divisible by $4$.

 $u=5$ has $4$ solutions:  every number ending in $5$ is divisible by $5$.

 $u=6$ has $1$ solution:  $36$.  $\pm 10$ or $\pm 20$ will kill divisibility by $3$, and thus kill divisibility by $6$.

 $u=7$ has no solutions.  The first multiples of $7$ that end in $7$ are $7$ and $77$, but both are outside of the range of this problem.

 $u=8$ has $1$ solution:  $48$.  $\pm 10, \pm 20, \pm 30$ will all kill divisibility by $8$ since $10, 20,$ and $30$ are not divisible by $8$.

 $u=9$ has no solutions.  $9$ and $99$ are the smallest multiples of $9$ that end in $9$.

 Totalling the solutions, we have $0 + 4 + 4 + 1 + 2 + 4 + 1 + 0 + 1 + 0 = 17$ solutions, giving the answer $\boxed{C}$, which is 17.

 


	\item (\textit{AMC 8 2000}) A block wall 100 feet long and 7 feet high will be constructed using blocks that are 1 foot high and either 2 feet long or 1 foot long (no blocks may be cut). The vertical joins in the blocks must be staggered as shown, and the wall must be even on the ends. What is the smallest number of blocks needed to build this wall?


 
                    \begin{figure}[H]
                        \centering
                        \includegraphics[width=0.5\textwidth]{Problem/AMC_8_AMC_2000_12_Problem_Figure_1.png}
                    \end{figure}
                    


 $\text{(A)}\ 344 \qquad \text{(B)}\ 347 \qquad \text{(C)}\ 350 \qquad \text{(D)}\ 353 \qquad \text{(E)}\ 356$


 
 \textbf{Answer: }D



\textbf{Solution}

Since the bricks are $1$ foot high, there will be $7$ rows.  To minimize the number of blocks used, rows $1, 3, 5,$ and $7$ will look like the bottom row of the picture, which takes $\frac{100}{2} = 50$ bricks to construct.  Rows $2, 4,$ and $6$ will look like the upper row pictured, which has $49$ 2-foot bricks in the middle, and $2$ 1-foot bricks on each end for a total of $51$ bricks.

 Four rows of $50$ bricks and three rows of $51$ bricks totals $4\cdot 50 + 3\cdot 51 = 200 + 153 = 353$ bricks, giving the answer $\boxed{D}.$

 If you wanted to be extra sure, you can realize that $700 / 2$ = 350 blocks. However, this is impossible, as you need them to be staggered.

 


	\item (\textit{AMC 8 2000}) In triangle $CAT$, we have $\angle ACT = \angle ATC$ and $\angle CAT = 36^\circ$. If $\overline{TR}$ bisects $\angle ATC$, then $\angle CRT =$


 
                    \begin{figure}[H]
                        \centering
                        \includegraphics[width=0.5\textwidth]{Problem/AMC_8_AMC_2000_13_Problem_Figure_1.png}
                    \end{figure}
                    


 $\text{(A)}\ 36^\circ \qquad \text{(B)}\ 54^\circ \qquad \text{(C)}\ 72^\circ \qquad \text{(D)}\ 90^\circ \qquad \text{(E)}\ 108^\circ$


 
 \textbf{Answer: }C



\textbf{Solution}

In $\triangle ACT$, the three angles sum to $180^\circ$, and $\angle C = \angle T$

 $\angle CAT + \angle ATC + \angle ACT = 180$

 $36 + \angle ATC + \angle ATC = 180$

 $2 \angle ATC = 144$

 $\angle ATC = 72$

 Since $\angle ATC$ is bisected by $\overline{TR}$, $\angle RTC = \frac{72}{2} = 36$

 Now focusing on the smaller $\triangle RTC$, the sum of the angles in that triangle is $180^\circ$, so:

 $\angle RTC + \angle TCR + \angle CRT = 180$

 $36 + \angle ACT + \angle CRT = 180$

 $36 + \angle ATC + \angle CRT = 180$

 $36 + 72 + \angle CRT = 180$

 $\angle CRT = 72^\circ$, giving the answer $\boxed{C}$

 


	\item (\textit{AMC 8 2000}) What is the units digit of $19^{19} + 99^{99}$?


 $\text{(A)}\ 0 \qquad \text{(B)}\ 1 \qquad \text{(C)}\ 2 \qquad \text{(D)}\ 8 \qquad \text{(E)}\ 9$


 
 \textbf{Answer: }D



\textbf{Solution}

Finding a pattern for each half of the sum, even powers of $19$ have a units digit of $1$, and odd powers of $19$ have a units digit of $9$. So, $19^{19}$ has a units digit of $9$. 

 Powers of $99$ have the exact same property, so $99^{99}$ also has a units digit of $9$. $9+9=18$ which has a units digit of $8$, so the answer is $\boxed{D}$.

 



\textbf{Solution 2}

Using modular arithmetic:
 \[99 \equiv 9 \equiv -1 \pmod{10}\]

 Similarly,
 \[19 \equiv 9 \equiv -1 \pmod{10}\]

 We have 
 \[(-1)^{19} + (-1)^{99} = -1 + -1 \equiv \boxed{(\textbf{D}) \ 8} \pmod{10}\]

 -ryjs

 



\textbf{Solution 3}

Experimentation gives
 \[\text{any number ending with }9^{\text{something even}} = \text{has units digit }1\]

 
 \[\text{any number ending with }9^{\text{something odd}} = \text{has units digit }9\]

 Using this we have
 \[19^{19} + 99^{99}\]
 \[9^{19} + 9^{99}\]

 Both $19$ and $99$ are odd, so we are left with
 \[9+9=18,\] which has units digit $\boxed{(\textbf{D}) \ 8}.$-ryjs

 



\textbf{Solution 4}

$19^{19}+99^{99}=$369729637649726772657187905628805440595668764281741102430259972423552570455277523421410650010128232727940978889548326540119429996769494359451621570193644014418071060667659303363419435659472789623878 which ends in 8. :) If you are really stupid and have time (like about 3 hours), this is a real brainless bash.

 --hefei417

 


	\item (\textit{AMC 8 2000}) Triangles $ABC$, $ADE$, and $EFG$ are all equilateral. Points $D$ and $G$ are midpoints of $\overline{AC}$ and $\overline{AE}$, respectively. If $AB = 4$, what is the perimeter of figure $ABCDEFG$?


 
                    \begin{figure}[H]
                        \centering
                        \includegraphics[width=0.5\textwidth]{Problem/AMC_8_AMC_2000_15_Problem_Figure_1.png}
                    \end{figure}
                    


 $\text{(A)}\ 12 \qquad \text{(B)}\ 13 \qquad \text{(C)}\ 15 \qquad \text{(D)}\ 18 \qquad \text{(E)}\ 21$


 
 \textbf{Answer: }C



\textbf{Solution 1}

The large triangle $ABC$ has sides of length $4$.  The medium triangle has sides of length $2$.  The small triangle has sides of length $1$.  There are $3$ segment sizes, and all segments depicted are one of these lengths.

 Starting at $A$ and going clockwise, the perimeter is:

 $AB + BC + CD + DE + EF + FG + GA$

 $4 + 4 + 2 + 2 + 1 + 1 + 1$

 $15$, thus the answer is $\boxed{C}$

 



\textbf{Solution 2}

The perimeter of $ABCDEFG$ is the perimeter of the three triangles, minus segments $AD$ and $EG$, which are on the interior of the figure.  Because each of these segments is on two triangles, each segment must be subtracted two times.

 As in solution 1, the sides of the triangles are $4, 2,$ and $1$, and the perimeters of the triangles are thus $12, 6,$ and $3$.

 The perimeter of the three triangles is $12 + 6 + 3 = 21$.  Subtracting the two segments $AD$ and $EG$ two times, the perimeter of $ABCDEFG$ is $21 - 2 - 1 - 2 - 1 = 15$, and the answer is $\boxed{C}$.

 



\textbf{Solution 3}

First, remove segments $DA$, and $EG$. Notice that segments $DE$, and $AG$ can fit into $DA$, and $EG$, respectively. This means we only need to account for the perimeters of triangles $ABC$, and $EFG$, whose perimeters are $12$, and $3$ respectively. Adding these perimeters gives us $12+3=15$, or answer choice $\boxed{C}$.~Electric\_Cheetah\_88

 


	\item (\textit{AMC 8 2000}) In order for Mateen to walk a kilometer (1000m) in his rectangular backyard, he must walk the length 25 times or walk its perimeter 10 times. What is the area of Mateen's backyard in square meters?


 $\text{(A)}\ 40 \qquad \text{(B)}\ 200 \qquad \text{(C)}\ 400 \qquad \text{(D)}\ 500 \qquad \text{(E)}\ 1000$


 
 \textbf{Answer: }C



\textbf{Solution}

The length $L$ of the rectangle is $\frac{1000}{25}=40$ meters. The perimeter $P$ is $\frac{1000}{10}=100$ meters. Since $P_{rect} = 2L + 2W$, we plug values in to get:

 $100 = 2\cdot 40 + 2W$

 $100 = 80 + 2W$

 $2W = 20$

 $W = 10$ meters

 Since $A_{rect} = LW$, the area is $40\cdot 10=400$ square meters or $\boxed{C}$.

 


	\item (\textit{AMC 8 2000}) The operation $\otimes$ is defined for all nonzero numbers by $a\otimes b = \dfrac{a^2}{b}$. Determine $[(1\otimes 2)\otimes 3] - [1\otimes (2\otimes 3)]$.


 $\text{(A)}\ -\dfrac{2}{3} \qquad \text{(B)}\ -\dfrac{1}{4} \qquad \text{(C)}\ 0 \qquad \text{(D)}\ \dfrac{1}{4} \qquad \text{(E)}\ \dfrac{2}{3}$


 
 \textbf{Answer: }A



\textbf{Solution}

Follow PE(MD)(AS), doing the innermost parentheses first.

 $[(1\otimes 2)\otimes 3]-[1\otimes (2\otimes 3)]$

 $[\frac{1^2}{2}\otimes 3]-[1\otimes \frac{2^2}{3}]$

 $[\frac{1}{2}\otimes 3]-[1\otimes \frac{4}{3}]$

 $[\frac{(\frac{1}{2})^{2}}{3}]-[\frac{1^2}{(\frac{4}{3})}]$

 $[\frac{1}{4} \cdot \frac{1}{3}]-[\frac{3}{4}]$

 $\frac{1}{12} - \frac{3}{4}$

 $\frac{1}{12} - \frac{9}{12}$

 $\frac{-8}{12}$

 $-\frac{2}{3}$, which is answer $\boxed{A}$

 


 


	\item (\textit{AMC 8 2000}) Consider these two geoboard quadrilaterals. Which of the following statements is true?


 
                    \begin{figure}[H]
                        \centering
                        \includegraphics[width=0.5\textwidth]{Problem/AMC_8_AMC_2000_18_Problem_Figure_1.png}
                    \end{figure}
                    


 $\text{(A)}\ \text{The area of quadrilateral I is more than the area of quadrilateral II.}$


 $\text{(B)}\ \text{The area of quadrilateral I is less than the area of quadrilateral II.}$


 $\text{(C)}\ \text{The quadrilaterals have the same area and the same perimeter.}$


 $\text{(D)}\ \text{The quadrilaterals have the same area, but the perimeter of I is more than the perimeter of II.}$


 $\text{(E)}\ \text{The quadrilaterals have the same area, but the perimeter of I is less than the perimeter of II.}$


 
 \textbf{Answer: }E



\textbf{Solution 1}

First consider the area of the two figures.  Assume that the pegs are $1$ unit apart.  Divide region I into two triangles by drawing a horizontal line on the second row from the top.  Shifting the bottom triangle up one unit will create a square, and this square has the same area as region I.  Thus, region I has an area of $1$ square unit.

 Draw a horizontal line on figure II on the fourth row from the top to divide the figure into two triangles.  Temporarily move the top of the higher triangle one peg to the left.  This will preserve the area of the triangle, as you are keeping both the height and the base of the triangle equal.  Now you have two congruent triangles that are half of the area of a unit square.  Thus, region II also has an area of $1$ square unit, and the areas are equal.

 To compare the perimeters, note that region I has two different side lengths: two sides are $1$ unit apart, and the other two sides are $\sqrt{2}$ apart, as they are the diagonal of a unit square.  The total perimeter is $2 + 2\sqrt{2}$.

 Note that region II has three different side lengths.  One side is a unit length, while two sides are $\sqrt{2}$.  For the perimeters to be equal, the last side must be of unit length.  But the last side in region II is clearly longer than $1$ unit, so we can say that the perimeter of region II is greater than the perimeter of region I without calculating it.

 Thus, the correct answer is $\boxed{E}$.

 (The perimeter of region II is $1 + 2\sqrt{2} + \sqrt{5}$, since the last side is the diagonal of a 2 by 1 rectangle, and can be found with the pythagorean theorem as $\sqrt{2^2 + 1^2} = \sqrt{5}$.  This does not need to be found for this problem, as you can do a one-to-one correspondence with three of the four sides of the figure as outlined, and just compare the last remaining side of each figure.)

 


 



\textbf{Solution 2}

Using the Pick’s Theorem, The area of both figures is 1. This rules out answer choices A and B. When using the Pythagorean Theorem, we find that the perimeter of I is 2 + 2($\sqrt{2}$) and the perimeter of II is 1 + 2($\sqrt{2}$) + $\sqrt{5}$. Since $\sqrt{5}$ > 1, the perimeter of II is greater, therefore the correct answer is $\boxed{E}$. Yaay!

 


	\item (\textit{AMC 8 2000}) Three circular arcs of radius 5 units bound the region shown. Arcs $AB$ and $AD$ are quarter-circles, and arc $BCD$ is a semicircle. What is the area, in square units, of the region?


 
                    \begin{figure}[H]
                        \centering
                        \includegraphics[width=0.5\textwidth]{Problem/AMC_8_AMC_2000_19_Problem_Figure_1.png}
                    \end{figure}
                    


 $\text{(A)}\ 25 \qquad \text{(B)}\ 10 + 5\pi \qquad \text{(C)}\ 50 \qquad \text{(D)}\ 50 + 5\pi \qquad \text{(E)}\ 25\pi$


 
 \textbf{Answer: }C



\textbf{Solution 1}

Draw two squares:  one that has opposing corners at $A$ and $B$, and one that has opposing corners at $A$ and $D$.  These squares share side $\overline{AO}$, where $O$ is the center of the large semicircle.

 These two squares have a total area of $2 \cdot 5^2$, but have two quarter circle "bites" of radius $5$ that must be removed.  Thus, the bottom part of the figure has area 

 $2\cdot 25 - 2 \cdot \frac{1}{4}\pi \cdot 5^2$

 $50 - \frac{25\pi}{2}$

 This is the area of the part of the figure underneath $\overline{BD}$.  The part of the figure over $\overline{BD}$ is just a semicircle with radius $5$, which has area of $\frac{1}{2}\pi\cdot 5^2 = \frac{25\pi}{2}$

 Adding the two areas gives a total area of $\boxed{(c)50}$

 



\textbf{Solution 2}

Draw line $\overline{BD}$.  Then draw $\overline {CO}$, where $O$ is the center of the semicircle.  You have two quarter circles on top, and two quarter circle-sized "bites" on the bottom.  Move the pieces from the top to fit in the bottom like a jigsaw puzzle.  You now have a rectangle with length $\overline {BD}$ and height $\overline {AO}$, which are equal to $10$ and $5$, respectively.  Thus, the total area is $50$, and the answer is $\boxed{C}$.

 


	\item (\textit{AMC 8 2000}) You have nine coins: a collection of pennies, nickels, dimes, and quarters having a total value of $$1.02$, with at least one coin of each type. How many dimes must you have?


 $\text{(A)}\ 1 \qquad \text{(B)}\ 2 \qquad \text{(C)}\ 3 \qquad \text{(D)}\ 4 \qquad \text{(E)}\ 5$


 
 \textbf{Answer: }A



\textbf{Solution 1}

Since you have one coin of each type, $1 + 5 + 10 + 25 = 41$ cents are already determined, leaving you with a total of $102 - 41 = 61$ cents remaining for $5$ coins.

 You must have $1$ more penny.  If you had more than $1$ penny, you must have at least $6$ pennies to leave a multiple of $5$ for the nickels, dimes, and quarters.  But you only have $5$ more coins to assign.

 Now you have $61 - 1 = 60$ cents remaining for $4$ coins, which may be nickels, quarters, or dimes.  If you have only one more dime, that leaves $50$ cents in $3$ nickels or quarters, which is impossible.  If you have two dimes, that leaves $40$ cents for $2$ nickels or quarters, which is again impossible.  If you have three dimes, that leaves $30$ cents for $1$ nickel or quarter, which is still impossible.  And all four remaining coins being dimes will not be enough.

 Therefore, you must have no more dimes to assign, and the $60$ cents in $4$ coins must be divided between the quarters and nickels.  We quickly see that $2$ nickels and $2$ quarters work.  Thus, the total count is $2$ quarters, $2$ nickels, $1$ penny, plus one more coin of each type that we originally subtracted.  Double-checking, that gives a total $2 + 2 + 1 + 4 = 9$ coins, and a total of $2\cdot 25 + 2\cdot 5 + 1 + (1 + 5 + 10 + 25) = 102$ cents.

 There is only $1$ dime in that combo, so the answer is $\boxed{A}$.

 


 



\textbf{Solution 2 (Faster)}

We see that there must be 102 cents, so therefore there's at least 2 pennies. That leaves 7 coins. We assume that there are 3 quarters, leaving 25 cents with 4 coins left. If all 4 are nickels, that would only be 20 cents, missing 5. Therefore, one nickel must be changed into 1 dime, so the answer is $\boxed{A}$

 Solution by \href{https://artofproblemsolving.com/wiki/index.php?title=User:ILoveMath31415926535}{ILoveMath31415926535}

 


 



\textbf{Solution 3}

It is clear that there should only be $2$ pennies; any more would take up too many coins, and the limit is $9$. Now we have $$1$ left, and $7$ coins to use. Looking at the quarters, we can make methodical guesses. If there is $1$ quarter, then we would have to make $$0.75$ with $6$ coins. We take a few educated guesses for the nickel and dime combos, and see that $1$ quarter will not work. Trying values for $2$ quarters, we see that this will not work either. When we reach $3$ quarters, the remaining is $$0.25$ made from $4$ coins. We try with $2$ dimes, which does not work (it only takes $3$ coins) and we try with $1$ dime. After trying $2$ pennies, $3$ quarters, $1$ dime, and $3$ nickels, it is evident that this combo works. Therefore, the answer is $\boxed{A}$

 
                \begin{itemize}
                    \item It might seem that this solution takes time, but in reality, it can be done very quickly, given ones desire to solve a problem quickly
                \end{itemize}


                


	\item (\textit{AMC 8 2000}) Keiko tosses one penny and Ephraim tosses two pennies. The probability that Ephraim gets the same number of heads that Keiko gets is:


 $\text{(A)}\ \dfrac{1}{4} \qquad \text{(B)}\ \dfrac{3}{8} \qquad \text{(C)}\ \dfrac{1}{2} \qquad \text{(D)}\ \dfrac{2}{3} \qquad \text{(E)}\ \dfrac{3}{4}$


 
 \textbf{Answer: }B



\textbf{Solution}

Divide it into $2$ cases:

 1) Keiko and Ephriam both get $0$ heads:This means that they both roll all tails, so there is only $1$ way for this to happen.

 2) Keiko and Ephriam both get $1$ head:For Keiko, there is only $1$ way for this to happen because he is only flipping 1 penny, but for Ephriam, there are 2 ways since there are $2$ choices for when he can flip the head. So, in total there are $2 \cdot 1 = 2$ ways for this case.

 Thus, in total there are $3$ ways that work. Since there are $2$ choices for each coin flip (Heads or Tails), there are $2^3 = 8$ total ways of flipping 3 coins.

 Thus, since all possible coin flips of 3 coins are equally likely, the probability is $\boxed{(B) \frac38}$.

 



\textbf{Solution 2}

Let $K(n)$ be the probability that Keiko gets $n$ heads, and let $E(n)$ be the probability that Ephriam gets $n$ heads.

 $K(0) = \frac{1}{2}$

 $K(1) = \frac{1}{2}$

 $K(2) = 0$  (Keiko only has one penny!)

 $E(0) = \frac{1}{2}\cdot\frac{1}{2} = \frac{1}{4}$

 $E(1) = \frac{1}{2}\cdot\frac{1}{2} + \frac{1}{2}\cdot\frac{1}{2} = 2\cdot\frac{1}{4} = \frac{1}{2}$  (because Ephraim can get HT or TH)

 $E(2) = \frac{1}{2}\cdot\frac{1}{2} = \frac{1}{4}$

 The probability that Keiko gets $0$ heads and Ephriam gets $0$ heads is $K(0)\cdot E(0)$.  Similarly for $1$ head and $2$ heads.  Thus, we have:

 $P = K(0)\cdot E(0) + K(1)\cdot E(1) + K(2)\cdot E(2)$

 $P = \frac{1}{2}\cdot\frac{1}{4} + \frac{1}{2}\cdot\frac{1}{2} + 0\cdot\frac{1}{4}$

 $P = \frac{3}{8}$

 Thus the answer is $\boxed{B}$.

 



\textbf{Solution 3}

Keiko can only get 1 head or one tail. That means Ephriam has to also get 1 head or one tail. There are 3 cases $HHT, HTH, TTT$ and each have a $\frac{1}{2}\cdot\frac{1}{2}\cdot\frac{1}{2}=\frac{1}{8}$ chance. Adding all 3 ways makes a total of $\boxed{(B)\frac38}$

 


 


	\item (\textit{AMC 8 2000}) A cube has edge length 2. Suppose that we glue a cube of edge length 1 on top of the big cube so that one of its faces rests entirely on the top face of the larger cube. The percent increase in the surface area (sides, top, and bottom) from the original cube to the new solid formed is closest to


 
                    \begin{figure}[H]
                        \centering
                        \includegraphics[width=0.5\textwidth]{Problem/AMC_8_AMC_2000_22_Problem_Figure_1.png}
                    \end{figure}
                    


 $\text{(A)}\ 10 \qquad \text{(B)}\ 15 \qquad \text{(C)}\ 17 \qquad \text{(D)}\ 21 \qquad \text{(E)}\ 25$


 
 \textbf{Answer: }C



\textbf{Solution 1}

The original cube has $6$ faces, each with an area of $2\cdot 2 = 4$ square units.  Thus the original figure had a total surface area of $24$ square units.

 The new figure has the original surface, with $6$ new faces that each have an area of $1$ square unit, for a total surface area of $6$ additional square units added to it.  But $1$ square unit of the top of the bigger cube, and $1$ square unit on the bottom of smaller cube, is not on the surface, and does not count towards the surface area.

 The total surface area is therefore $24 + 6 - 1 - 1 = 28$ square units.

 The percent increase in surface area is $\frac{SA_{new} - SA_{old}}{SA_{old}}\cdot 100\% = \frac{28-24}{24}\cdot 100\% \approx 16.67\%$, giving the closest answer as $\boxed{C}$.

 



\textbf{Solution 2(Similar to 1)}

The original cube has $6$ faces, each with an area of $2\cdot 2 = 4$ square units.  Thus the original figure had a total surface area of $4$ times $6$ square units. Notice that the little cube covers a bit of the larger cube but this is replaced by the top of the little cube. Therefore, the actual increase comes from the four sides of the cube giving $1$ squared = $1$, $1$ ($4$) = $4$. So $4$ over the original $4$ times $6$. 4/4(6) = 1/6which is about 16.67\% giving $\boxed{C}$..

 



\textbf{Solution 3}

To be honest, we could just bash this one out to get about 20\%, which is C. Yaaaay! BRO, who ever made this is tweaking, if you bashed it to 20 you would choose D so this is not bashable. Tsk Tsk  Plus they never said how to bash¯\_()D is closer to 20 sir lock in

 


	\item (\textit{AMC 8 2000}) There is a list of seven numbers. The average of the first four numbers is 5, and the average of the last four numbers is 8. If the average of all seven numbers is $6\frac{4}{7}$, then the number common to both sets of four numbers is


 $\text{(A)}\ 5\frac{3}{7} \qquad \text{(B)}\ 6 \qquad \text{(C)}\ 6\frac{4}{7} \qquad \text{(D)}\ 7 \qquad \text{(E)}\ 7\frac{3}{7}$


 
 \textbf{Answer: }B



\textbf{Solution}

Remember that if a list of $n$ numbers has an average of $k$, then the sum $S$ of all the numbers on the list is $S = nk$.

 So if the average of the first $4$ numbers is $5$, then the first four numbers total $4 \cdot 5 = 20$.

 If the average of the last $4$ numbers is $8$, then the last four numbers total $4 \cdot 8 = 32$.

 If the average of all $7$ numbers is $6\frac{4}{7}$, then the total of all seven numbers is $7 \cdot 6\frac{4}{7} = 7\cdot 6 + 4 =  46$.

 If the first four numbers are $20$, and the last four numbers are $32$, then all "eight" numbers are $20 + 32 = 52$.  But that's counting one number twice.  Since the sum of all seven numbers is $46$, then the number that was counted twice is $52 - 46 = 6$, and the answer is $\boxed{(B) 6}$

 Algebraically, if $a + b + c + d = 20$, and $d + e + f + g = 32$, you can add both equations to get $a + b + c + 2d + e + f + g = 52$.  You know that $a + b + c + d + e + f + g = 46$, so you can subtract that from the last equation to get $d = \boxed{(B) 6}$, and $d$ is the number that appeared twice.

 Yay!  :D

 


 


	\item (\textit{AMC 8 2000}) If $\angle A = 20^\circ$ and $\angle AFG = \angle AGF$, then $\angle B + \angle D =$


 
                    \begin{figure}[H]
                        \centering
                        \includegraphics[width=0.5\textwidth]{Problem/AMC_8_AMC_2000_24_Problem_Figure_1.png}
                    \end{figure}
                    


 $\text{(A)}\ 48^\circ \qquad \text{(B)}\ 60^\circ \qquad \text{(C)}\ 72^\circ \qquad \text{(D)}\ 80^\circ \qquad \text{(E)}\ 90^\circ$


 
 \textbf{Answer: }D



\textbf{Solution}

As a strategy, think of how $\angle B + \angle D$ would be determined, particularly without determining either of the angles individually, since it may not be possible to determine $\angle B$ or $\angle D$ alone.  If you see $\triangle BFD$, then you can see that the problem is solved quickly after determining $\angle BFD$.

 But start with $\triangle AGF$, since that's where most of our information is.  Looking at $\triangle AGF$, since $\angle AFG = \angle AGF$, and $\angle A = 20$, we can write:

 $\angle A + \angle AFG + \angle AGF = 180$

 $20 + 2\angle AFG = 180$

 $\angle AFG = 80$

 By noting that $\angle AFG$ and $\angle GFD$ make a straight line, we know

 $\angle AFG + \angle GFD = 180$

 $80 + \angle GFD = 180$

 $\angle GFD = 100$

 Ignoring all other parts of the figure and looking only at $\triangle BFD$, you see that $\angle B + \angle D + \angle BFD = 180$.  But $\angle BFD$ is the same as $\angle GFD$.  Therefore:

 $\angle B + \angle D + \angle GFD = 180$

 $\angle B + \angle D + 100 = 180$

 $\angle B + \angle D = 80^\circ$, and the answer is thus $\boxed{D}$

 


	\item (\textit{AMC 8 2000}) The area of rectangle $ABCD$ is 72. If point $A$ and the midpoints of $\overline{BC}$ and $\overline{CD}$ are joined to form a triangle, the area of that triangle is


 
                    \begin{figure}[H]
                        \centering
                        \includegraphics[width=0.5\textwidth]{Problem/AMC_8_AMC_2000_25_Problem_Figure_1.png}
                    \end{figure}
                    


 $\text{(A)}\ 21 \qquad \text{(B)}\ 27 \qquad \text{(C)}\ 30 \qquad \text{(D)}\ 36 \qquad \text{(E)}\ 40$


 
 \textbf{Answer: }B



\textbf{Solution 1}

To quickly solve this multiple choice problem, make the (not necessarily valid, but very convenient) assumption that $ABCD$ can have any dimension.  Give the rectangle dimensions of $AB = CD = 12$ and $BC = AD= 6$, which is the easiest way to avoid fractions.  Labeling the right midpoint as $M$, and the bottom midpoint as $N$, we know that $DN = NC = 6$, and $BM = MC = 3$.  

 $[\triangle ADN] = \frac{1}{2}\cdot 6\cdot 6 = 18$

 $[\triangle MNC] = \frac{1}{2}\cdot 3\cdot 6 = 9$

 $[\triangle ABM] = \frac{1}{2}\cdot 12\cdot 3 = 18$

 $[\triangle AMN] = [\square ABCD] - [\triangle ADN] - [\triangle MNC] - [\triangle ABM]$

 $[\triangle AMN] = 72 - 18 - 9 - 18$

 $[\triangle AMN] = 27$, and the answer is $\boxed{B}$ 

 



\textbf{Solution 2}

The above answer is fast, but satisfying, and assumes that the area of $\triangle AMN$ is independent of the dimensions of the rectangle. All in all, it's a very good answer though. However this is an alternative if you don't get the above answer. Label $AB = CD = l$ and $BC = DA = h$

 Labelling $m$ and $n$ as the right and lower midpoints respectively, and redoing all the work above, we get:

 $[\triangle ADN] = \frac{1}{2}\cdot h\cdot \frac{l}{2} = \frac{lh}{4}$

 $[\triangle MNC] = \frac{1}{2}\cdot \frac{l}{2}\cdot \frac{h}{2} = \frac{lh}{8}$

 $[\triangle ABM] = \frac{1}{2}\cdot l\cdot \frac{h}{2} = \frac{lh}{4}$

 $[\triangle AMN] = [\square ABCD] - [\triangle ADN] - [\triangle MNC] - [\triangle ABM]$

 $[\triangle AMN] = lh - \frac{lh}{4} - \frac{lh}{8} - \frac{lh}{4}$

 $[\triangle AMN] = \frac{3}{8}lh = \frac{3}{8}\cdot 72 = 27$, and the answer is $\boxed{B}$

 



\textbf{Solution 3}

Let's assume WLOG that the sides of the rectangle are $9$ and $8.$ The area of the 3 triangles would then be $8\cdot\frac{9}{2}\cdot\frac{1}{2} = 18,$$4\cdot\frac{9}{2}\cdot\frac{1}{2} = 9,$$4\cdot 9\cdot\frac{1}{2} = 18.$Adding these up, we get $45$, and subtracting that from $72$, we get $27$, so the answer is $\boxed{B}$

 


\end{enumerate}